\chapter{Datenvorbereitung und Evaluation}
\label{chap:datenvorbereitung}

\begin{myremark}[Algorithmusart in dieser DL]
Funktionsart: \textbf{Methodik- und Evaluationsbaustein} für Klassifikation, Filterung und andere datengetriebene Verfahren. \\
Umsetzungsart: \textbf{algorithmusübergreifend}, kein eigener Vorhersagealgorithmus.
\end{myremark}

\begin{myremark}
Bevor ein Modell trainiert werden kann, müssen die Daten vorbereitet werden.
Ausserdem brauchen wir Werkzeuge, um die Qualität eines Modells zu messen.
\end{myremark}

\section{Trainings- und Testdaten}

\begin{mydefinition}[Train-Test-Split]
Die Daten werden in einen \textbf{Trainings\-datensatz} und einen
\textbf{Testdatensatz} aufgeteilt (typisch: 80\,\% / 20\,\%).
Das Modell wird \emph{nur} mit den Trainingsdaten trainiert und
anschliessend auf den \emph{unseen} Testdaten bewertet.
\end{mydefinition}

\begin{mycaution}
Wer das Modell auch auf den Trainingsdaten bewertet, misst lediglich,
wie gut das Modell \enquote{auswendig gelernt} hat –
das gibt keinen Aufschluss über die echte Leistung auf neuen Daten.
\end{mycaution}

\begin{myexercise}[Train-Test-Split nachvollziehen]
Ein Datensatz enthält 500 Zeilen.
\begin{enumerate}
    \item Wie viele Zeilen enthält der Trainings- bzw. der Testdatensatz bei einem 80/20-Split?
    \item Warum darf man die Testdaten \emph{nicht} für die Modellentwicklung verwenden?
    \item Was ist \emph{Datenleckage} (data leakage) und warum ist sie problematisch?
\end{enumerate}
\begin{myanswer}
1.\ 400 Training / 100 Test \quad
2.\ Sonst misst man nur, ob das Modell die Testdaten auswendig kennt, nicht ob es generalisiert. \quad
3.\ Information aus dem Testset fliesst in das Training – der gemessene Fehler unterschätzt den echten Fehler.
\end{myanswer}
\end{myexercise}

\section{Merkmale und Zielvariable}

\begin{mydefinition}[Features und Label]
\begin{itemize}
    \item \textbf{Features} (Merkmale): die Eingabevariablen $x_1, x_2, \ldots, x_n$
    \item \textbf{Label} (Zielvariable): der Ausgabewert $y$, den das Modell vorhersagen soll
\end{itemize}
\end{mydefinition}

\begin{myexample}
Beim Datensatz \texttt{learning\_style\_classifier.csv}:
\begin{itemize}
    \item Features: Lernzeit, Schlafstunden, Anzahl Wiederholungen
    \item Label: Lernstil (\enquote{visuell}, \enquote{auditiv}, \enquote{kinästhetisch})
\end{itemize}
\end{myexample}

\section{Normalisierung}

\begin{mydefinition}[Min-Max-Normalisierung]
Features mit sehr unterschiedlichen Wertebereichen können Modelle verzerren.
Die \textbf{Min-Max-Normalisierung} skaliert alle Werte auf $[0,\,1]$:
\[
x' = \frac{x - x_{\min}}{x_{\max} - x_{\min}}
\]
\end{mydefinition}

\begin{myexercise}[Normalisierung berechnen]
Gegeben: Lernzeiten $\{2, 4, 6, 8, 10\}$ Stunden.
\begin{enumerate}
    \item Normalisieren Sie alle Werte mit Min-Max.
    \item Welcher Wert entspricht nach der Normalisierung dem Wert 6?
\end{enumerate}
\begin{myanswer}
$x_{\min}=2,\ x_{\max}=10$: 
$\{0,\; 0.25,\; 0.5,\; 0.75,\; 1.0\}$. 
Der Wert 6 wird zu $0.5$.
\end{myanswer}
\end{myexercise}

\section{Gütemasse für Klassifikation}

\begin{mydefinition}[Konfusionsmatrix]
Die \textbf{Konfusionsmatrix} zeigt, wie oft das Modell richtig und falsch klassifiziert hat:

\begin{center}
\begin{tabular}{cc|cc}
 & & \multicolumn{2}{c}{\textbf{Vorhergesagt}} \\
 & & Positiv & Negativ \\
\hline
\multirow{2}{*}{\textbf{Tatsächlich}} & Positiv & TP & FN \\
 & Negativ & FP & TN \\
\end{tabular}
\end{center}
\end{mydefinition}

\begin{mydefinition}[Accuracy, Precision, Recall]
\[
\text{Accuracy} = \frac{TP + TN}{TP + TN + FP + FN}
\qquad
\text{Precision} = \frac{TP}{TP + FP}
\qquad
\text{Recall} = \frac{TP}{TP + FN}
\]

Die \textbf{Accuracy} misst die Gesamtgenauigkeit des Modells.
Die \textbf{Precision} misst, wie viele der als positiv vorhergesagten Fälle tatsächlich positiv sind.
Der \textbf{Recall} misst, wie viele der tatsächlich positiven Fälle korrekt erkannt wurden.

Die Begriffe ``Sensitivität'' und ``Spezifizität'' werden manchmal synonym zu Recall bzw. Precision verwendet.
\end{mydefinition}

\begin{myexercise}[Metriken berechnen]
Ein Spam-Filter klassifiziert 200 E-Mails.
Die Konfusionsmatrix ergibt: 
\begin{center}
\begin{tabular}{cc|cc}
 & & \multicolumn{2}{c}{\textbf{Vorhergesagt}} \\
 & & Spam & Kein Spam \\
\hline
\multirow{2}{*}{\textbf{Tatsächlich}} & Spam & 80 & 20 \\
 & Kein Spam & 10 & 90 \\
\end{tabular}
\end{center}

\begin{enumerate}
    \item Berechnen Sie Accuracy, Precision und Recall.
    \item Warum ist Recall besonders wichtig, wenn es darum geht, keinen Spam zu verpassen?
\end{enumerate}
\begin{myanswer}
Accuracy $= \frac{80+90}{200} = 0.85$; \quad
Precision $= \frac{80}{90} \approx 0.89$; \quad
Recall $= \frac{80}{100} = 0.80$. \\
Ein hoher Recall stellt sicher, dass möglichst wenig Spam als \enquote{kein Spam} durchrutscht (wenig FN).
\end{myanswer}
\end{myexercise}

\begin{myremark}[Notebook]
Praktische Umsetzung:
\href{\notebookbaseurl00_Datenvorbereitung.ipynb}{\texttt{00\_Datenvorbereitung.ipynb}}
\end{myremark}
